\begin{thebibliography}{99}
\footnotesize\itemsep=0pt

\bibitem{AM08}
Adamovi\'c D., Milas A., On the triplet vertex algebra {$\mathcal W(p)$},
 \href{https://doi.org/10.1016/j.aim.2007.11.012}{\textit{Adv. Math.}} \textbf{217} (2008), 2664--2699, \href{https://arxiv.org/abs/0707.1857}{arXiv:0707.1857}.

\bibitem{AM14}
Adamovi\'c D., Milas A., {$C_2$}-cofinite {$\mathcal W$}-algebras and their
 logarithmic representations, in Conformal Field Theories and Tensor
 Categories, \textit{Math. Lect. Peking Univ.}, \href{https://doi.org/10.1007/978-3-642-39383-9_6}{Springer}, Heidelberg, 2014, 249--270,
 \href{https://arxiv.org/abs/1212.6771}{arXiv:1212.6771}.

\bibitem{AA17}
Andruskiewitsch N., Angiono I., On finite dimensional {N}ichols algebras of
 diagonal type, \href{https://doi.org/10.1007/s13373-017-0113-x}{\textit{Bull. Math. Sci.}} \textbf{7} (2017), 353--573,
 \href{https://arxiv.org/abs/1707.08387}{arXiv:1707.08387}.

\bibitem{AAB14}
Andruskiewitsch N., Angiono I., Bagio D., Examples of pointed color {H}opf
 algebras, \href{https://doi.org/10.1142/S0219498813500989}{\textit{J.~Algebra Appl.}} \textbf{13} (2014), 1350098, 28~pages,
 \href{https://arxiv.org/abs/1212.0514}{arXiv:1212.0514}.

\bibitem{AAR19}
Andruskiewitsch N., Angiono I., Rossi~Bertone F., Lie algebras arising from
 {N}ichols algebras of diagonal type, \href{https://arxiv.org/abs/1911.06586}{arXiv:1911.06586}.

\bibitem{AAY11}
Andruskiewitsch N., Angiono I., Yamane H., On pointed {H}opf superalgebras, in
 New Developments in {L}ie Theory and Its Applications, \textit{Contemp.
 Math.}, Vol.~544, \href{https://doi.org/10.1090/conm/544/10752}{Amer. Math. Soc.}, Providence, RI, 2011, 123--140,
 \href{https://arxiv.org/abs/1009.5148}{arXiv:1009.5148}.

\bibitem{AHS10}
Andruskiewitsch N., Heckenberger I., Schneider H.-J., The {N}ichols algebra of a
 semisimple {Y}etter--{D}rinfeld module, \textit{Amer.~J. Math.} \textbf{132}
 (2010), 1493--1547, \href{https://arxiv.org/abs/0803.2430}{arXiv:0803.2430}.

\bibitem{AS10}
Andruskiewitsch N., Schneider H.-J., On the classification of finite-dimensional
 pointed {H}opf algebras, \href{https://doi.org/10.4007/annals.2010.171.375}{\textit{Ann. of Math.}} \textbf{171} (2010),
 375--417, \href{https://arxiv.org/abs/math.QA/0502157}{arXiv:math.QA/0502157}.

\bibitem{Ang13}
Angiono I., On {N}ichols algebras of diagonal type, \href{https://doi.org/10.1515/crelle-2011-0008}{\textit{J.~Reine Angew.
 Math.}} \textbf{683} (2013), 189--251, \href{https://arxiv.org/abs/1104.0268}{arXiv:1104.0268}.

\bibitem{Ang16}
Angiono I., Distinguished pre-{N}ichols algebras, \href{https://doi.org/10.1007/s00031-015-9341-x}{\textit{Transform. Groups}}
 \textbf{21} (2016), 1--33, \href{https://arxiv.org/abs/1405.6681}{arXiv:1405.6681}.

\bibitem{Arak07}
Arakawa T., Representation theory of {$\mathcal W$}-algebras, \href{https://doi.org/10.1007/s00222-007-0046-1}{\textit{Invent.
 Math.}} \textbf{169} (2007), 219--320, \href{https://arxiv.org/abs/math.QA/0506056}{arXiv:math.QA/0506056}.

\bibitem{CLR21}
Auger J., Creutzig T., Kanade S., Rupert M., Characterizing braided tensor
 categories associated to logarithmic vertex operator algebras,
 \href{https://arxiv.org/abs/2104.13262}{arXiv:2104.13262}.

\bibitem{BLS15}
Barvels A., Lentner S., Schweigert C., Partially dualized {H}opf algebras have
 equivalent {Y}etter--{D}rinfel'd modules, \href{https://doi.org/10.1016/j.jalgebra.2015.02.010}{\textit{J.~Algebra}} \textbf{430}
 (2015), 303--342, \href{https://arxiv.org/abs/1402.2214}{arXiv:1402.2214}.

\bibitem{CGR20}
Creutzig T., Gainutdinov A.M., Runkel I., A quasi-{H}opf algebra for the
 triplet vertex operator algebra, \href{https://doi.org/10.1142/S021919971950024X}{\textit{Commun. Contemp. Math.}} \textbf{22}
 (2020), 1950024, 71~pages, \href{https://arxiv.org/abs/1712.07260}{arXiv:1712.07260}.

\bibitem{Cun10}
Cuntz M., Crystallographic arrangements: {W}eyl groupoids and simplicial
 arrangements, \href{https://doi.org/10.1112/blms/bdr009}{\textit{Bull. Lond. Math. Soc.}} \textbf{43} (2011), 734--744,
 \href{https://arxiv.org/abs/1006.1997}{arXiv:1006.1997}.

\bibitem{CH09}
Cuntz M., Heckenberger I., Finite {W}eyl groupoids of rank three,
 \href{https://doi.org/10.1090/S0002-9947-2011-05368-7}{\textit{Trans. Amer. Math. Soc.}} \textbf{364} (2012), 1369--1393,
 \href{https://arxiv.org/abs/0912.0212}{arXiv:0912.0212}.

\bibitem{CH10}
Cuntz M., Heckenberger I., Finite {W}eyl groupoids, \href{https://doi.org/10.1515/crelle-2013-0033}{\textit{J.~Reine Angew.
 Math.}} \textbf{702} (2015), 77--108, \href{https://arxiv.org/abs/1008.5291}{arXiv:1008.5291}.

\bibitem{DL93}
Dong C., Lepowsky J., Generalized vertex algebras and relative vertex
 operators, \textit{Progress in Mathematics}, Vol.~112, \href{https://doi.org/10.1007/978-1-4612-0353-7}{Birkh\"auser Boston},
 Inc., Boston, MA, 1993.

\bibitem{DF84}
Dotsenko V.S., Fateev V.A., Conformal algebra and multipoint correlation
 functions in {$2$}{D} statistical models, \href{https://doi.org/10.1016/0550-3213(84)90269-4}{\textit{Nuclear Phys.~B}}
 \textbf{240} (1984), 312--348.

\bibitem{FF84}
Feigin B.L., Fuchs D.B., Verma modules over the {V}irasoro algebra, in Topology
 ({L}eningrad, 1982), \textit{Lecture Notes in Math.}, Vol.~1060, \href{https://doi.org/10.1007/BFb0099939}{Springer},
 Berlin, 1984, 230--245.

\bibitem{FT10}
Feigin B.L., Tipunin I.Yu., Logarithmic {CFT}s connected with simple {L}ie
 algebras, \href{https://arxiv.org/abs/1002.5047}{arXiv:1002.5047}.

\bibitem{Fel89}
Felder G., B{RST} approach to minimal models, \href{https://doi.org/10.1016/0550-3213(89)90568-3}{\textit{Nuclear Phys.~B}}
 \textbf{317} (1989), 215--236.

\bibitem{FL17}
Flandoli I., Lentner S., Logarithmic conformal field theories of type~{$B_n$},
 {$\ell=4$} and symplectic fermions, \href{https://doi.org/10.1063/1.5010904}{\textit{J.~Math. Phys.}} \textbf{59}
 (2018), 071701, 35~pages, \href{https://arxiv.org/abs/1706.07994}{arXiv:1706.07994}.

\bibitem{FSS96}
Frappat L., Sciarrino A., Sorba P., Dictionary on {L}ie superalgebras,
 \href{https://arxiv.org/abs/hep-th/9607161}{arXiv:hep-th/9607161}.

\bibitem{FBZ04}
Frenkel E., Ben-Zvi D., Vertex algebras and algebraic curves, 2nd ed.,
 \textit{Mathematical Surveys and Monographs}, Vol.~88, \href{https://doi.org/10.1090/surv/088}{Amer. Math.
 Soc.}, Providence, RI, 2004.

\bibitem{GRW09}
Gaberdiel M.R., Runkel I., Wood S., Fusion rules and boundary conditions in the
 {$c=0$} triplet model, \href{https://doi.org/10.1088/1751-8113/42/32/325403}{\textit{J.~Phys.~A: Math. Theor.}} \textbf{42} (2009),
 325403, 43~pages, \href{https://arxiv.org/abs/0905.0916}{arXiv:0905.0916}.

\bibitem{GLO18}
Gainutdinov A.M., Lentner S., Ohrmann T., Modularization of small quantum
 groups, \href{https://arxiv.org/abs/1809.02116}{arXiv:1809.02116}.

\bibitem{FGST06a}
Gainutdinov A.M., Semikhatov A.M., Tipunin I.Yu., Feigin B.L.,
 Kazhdan--{L}usztig correspondence for the representation category of the
 triplet $W$-algebra in logarithmic {CFT}, \href{https://doi.org/10.1007/s11232-006-0113-6}{\textit{Theoret. and Math. Phys.}}
 \textbf{148} (2006), 1210--1235, \href{https://arxiv.org/abs/math.QA/0512621}{arXiv:math.QA/0512621}.

\bibitem{G21}
Gaitsgory D., A conjectural extension of the {K}azhdan--{L}usztig equivalence,
 \href{https://doi.org/10.4171/prims/57-3-14}{\textit{Publ. Res. Inst. Math. Sci.}} \textbf{57} (2021), 1227--1376,
 \href{https://arxiv.org/abs/1810.09054}{arXiv:1810.09054}.

\bibitem{GN21}
Gannon T., Negron C., Quantum {${\rm SL}(2)$} and logarithmic vertex operator
 algebras at $(p,1)$-central charge, \href{https://arxiv.org/abs/2104.12821}{arXiv:2104.12821}.

\bibitem{GG17}
Garcia G.A., Gavarini F., Multiparameter quantum groups at roots of unity,
 \href{https://arxiv.org/abs/1708.05760}{arXiv:1708.05760}.

\bibitem{Heck05}
Heckenberger I., Weyl equivalence for rank 2 {N}ichols algebras of diagonal
 type, \href{https://doi.org/10.1007/BF02824835}{\textit{Ann. Univ. Ferrara}} \textbf{51} (2005), 281--289,
 \href{https://arxiv.org/abs/math.QA/0411621}{arXiv:math.QA/0411621}.

\bibitem{Hecklist}
Heckenberger I., Classification of arithmetic root systems of rank 3, in
 Proceedings of the {XVI}th {L}atin {A}merican {A}lgebra {C}olloquium
 ({S}panish), \textit{Bibl. Rev. Mat. Iberoamericana}, Rev. Mat. Iberoamericana,
 Madrid, 2007, 227--252, \href{https://arxiv.org/abs/math.QA/0509145}{arXiv:math.QA/0509145}.

\bibitem{Heck06}
Heckenberger I., Classification of arithmetic root systems, \href{https://doi.org/10.1016/j.aim.2008.08.005}{\textit{Adv. Math.}}
 \textbf{220} (2009), 59--124, \href{https://arxiv.org/abs/math.QA/0605795}{arXiv:math.QA/0605795}.

\bibitem{Heck07}
Heckenberger I., Lusztig isomorphisms for {D}rinfel'd doubles of bosonizations
 of {N}ichols algebras of diagonal type, \href{https://doi.org/10.1016/j.jalgebra.2010.02.013}{\textit{J.~Algebra}} \textbf{323}
 (2010), 2130--2182, \href{https://arxiv.org/abs/0710.4521}{arXiv:0710.4521}.

\bibitem{HS08}
Heckenberger I., Schneider H.-J., Root systems and {W}eyl groupoids for
 {N}ichols algebras, \href{https://doi.org/10.1112/plms/pdq001}{\textit{Proc. Lond. Math. Soc.}} \textbf{101} (2010),
 623--654, \href{https://arxiv.org/abs/0807.0691}{arXiv:0807.0691}.

\bibitem{HS11}
Heckenberger I., Schneider H.-J., Yetter--{D}rinfeld modules over bosonizations
 of dually paired {H}opf algebras, \href{https://doi.org/10.1016/j.aim.2013.05.009}{\textit{Adv. Math.}} \textbf{244} (2013),
 354--394, \href{https://arxiv.org/abs/1111.4673}{arXiv:1111.4673}.

\bibitem{HS20}
Heckenberger I., Schneider H.-J., Hopf algebras and root systems,
 \textit{Mathematical Surveys and Monographs}, Vol.~247, Amer. Math. Soc.,
 Providence, RI, 2020.

\bibitem{HY08}
Heckenberger I., Yamane H., A generalization of {C}oxeter groups, root systems,
 and {M}atsumoto's theorem, \href{https://doi.org/10.1007/s00209-007-0223-3}{\textit{Math.~Z.}} \textbf{259} (2008), 255--276,
 \href{https://arxiv.org/abs/math.QA/0610823}{arXiv:math.QA/0610823}.

\bibitem{Helb10}
Helbig M., On the lifting of {N}ichols algebras, \href{https://doi.org/10.1080/00927872.2011.568029}{\textit{Comm. Algebra}}
 \textbf{40} (2012), 3317--3351, \href{https://arxiv.org/abs/1003.5882}{arXiv:1003.5882}.

\bibitem{HLZ06}
Huang Y.-Z., Lepowsky J., Zhang L., A logarithmic generalization of tensor
 product theory for modules for a vertex operator algebra,
 \href{https://doi.org/10.1142/S0129167X06003758}{\textit{Internat.~J. Math.}} \textbf{17} (2006), 975--1012,
 \href{https://arxiv.org/abs/math.QA/0311235}{arXiv:math.QA/0311235}.

\bibitem{HLZ10}
Huang Y.-Z., Lepowsky J., Zhang L., Logarithmic tensor category theory for
 generalized modules for a conformal vertex algebra,~I: {I}ntroduction and
 strongly graded algebras and their generalized modules, \href{https://arxiv.org/abs/1012.4193}{arXiv:1012.4193}.

\bibitem{Kac77}
Kac V.G., Lie superalgebras, \href{https://doi.org/10.1016/0001-8708(77)90017-2}{\textit{Adv. Math.}} \textbf{26} (1977),
 8--96.

\bibitem{Kac98}
Kac V.G., Vertex algebras for beginners, 2nd ed., \textit{University Lecture Series},
 Vol.~10, \href{https://doi.org/10.1090/ulect/010}{Amer. Math. Soc.}, Providence, RI, 1998.

\bibitem{Kad97}
Kadell K.W.J., The {S}elberg--{J}ack symmetric functions, \href{https://doi.org/10.1006/aima.1997.1642}{\textit{Adv. Math.}}
 \textbf{130} (1997), 33--102.

\bibitem{KT91}
Khoroshkin S.M., Tolstoy V.N., Universal {$R$}-matrix for quantized
 (super)algebras, \href{https://doi.org/10.1007/BF02102819}{\textit{Comm. Math. Phys.}} \textbf{141} (1991), 599--617.

\bibitem{Len14}
Lentner S., New large-rank {N}ichols algebras over nonabelian groups with
 commutator subgroup {$\mathbb{Z}_2$}, \href{https://doi.org/10.1016/j.jalgebra.2014.07.017}{\textit{J.~Algebra}} \textbf{419}
 (2014), 1--33, \href{https://arxiv.org/abs/1306.5684}{arXiv:1306.5684}.

\bibitem{L19}
Lentner S., The unrolled quantum group inside {L}usztig's quantum group of
 divided powers, \href{https://doi.org/10.1007/s11005-019-01185-9}{\textit{Lett. Math. Phys.}} \textbf{109} (2019), 1665--1682,
 \href{https://arxiv.org/abs/1702.05164}{arXiv:1702.05164}.

\bibitem{Len17}
Lentner S., Quantum groups and {N}ichols algebras acting on conformal field
 theories, \href{https://doi.org/10.1016/j.aim.2020.107517}{\textit{Adv. Math.}} \textbf{378} (2021), 107517, 71~pages,
 \href{https://arxiv.org/abs/1702.06431}{arXiv:1702.06431}.

\bibitem{Lusz93}
Lusztig G., Introduction to quantum groups, \textit{Progress in Mathematics},
 Vol.~110, Birkh\"auser Boston, Inc., Boston, MA, 1993.

\bibitem{Mac95}
Macdonald I.G., Symmetric functions and {H}all polynomials, 2nd ed., \textit{Oxford
 Mathematical Monographs}, The Clarendon Press, Oxford University Press, New
 York, 1995.

\bibitem{MV00}
Mukhin E., Varchenko A., Remarks on critical points of phase functions and
 norms of {B}ethe vectors, in Arrangements~-- {T}okyo 1998, \textit{Adv. Stud.
 Pure Math.}, Vol.~27, \href{https://doi.org/10.2969/aspm/02710239}{Kinokuniya}, Tokyo, 2000, 239--246,
 \href{https://arxiv.org/abs/math.RT/9810087}{arXiv:math.RT/9810087}.

\bibitem{NT11}
Nagatomo K., Tsuchiya A., The triplet vertex operator algebra {$W(p)$} and the
 restricted quantum group {$\overline U_q(\mathfrak{sl}_2)$} at {$q={\rm
 e}^{\frac{\pi {\rm i}}{p}}$}, in Exploring new structures and natural constructions
 in mathematical physics, \textit{Adv. Stud. Pure Math.}, Vol.~61, \href{https://doi.org/10.2969/aspm/06110001}{Math. Soc.
 Japan}, Tokyo, 2011, 1--49, \href{https://arxiv.org/abs/0902.4607}{arXiv:0902.4607}.

\bibitem{Ne21}
Negron C., Log-modular quantum groups at even roots of unity and the quantum
 {F}robenius~{I}, \href{https://doi.org/10.1007/s00220-021-04012-2}{\textit{Comm. Math. Phys.}} \textbf{382} (2021), 773--814,
 \href{https://arxiv.org/abs/1812.02277}{arXiv:1812.02277}.

\bibitem{Rosso98}
Rosso M., Quantum groups and quantum shuffles, \href{https://doi.org/10.1007/s002220050249}{\textit{Invent. Math.}}
 \textbf{133} (1998), 399--416.

\bibitem{Sel44}
Selberg A., Remarks on a multiple integral, \textit{Norsk Mat. Tidsskr.}
 \textbf{26} (1944), 71--78.

\bibitem{Sem14}
Semikhatov A.M., Virasoro central charges for {N}ichols algebras, in Conformal
 field theories and tensor categories, \textit{Math. Lect. Peking Univ.}, \href{https://doi.org/10.1007/978-3-642-39383-9_3}{Springer},
 Heidelberg, 2014, 67--92, \href{https://arxiv.org/abs/1109.1767}{arXiv:1109.1767}.

\bibitem{ST13}
Semikhatov A.M., Tipunin I.Yu., Logarithmic {$\widehat{\mathfrak{sl}}(2)$} {CFT}
 models from {N}ichols algebras:~{I}, \href{https://doi.org/10.1088/1751-8113/46/49/494011}{\textit{J.~Phys.~A: Math. Theor.}}
 \textbf{46} (2013), 494011, 53~pages, \href{https://arxiv.org/abs/1301.2235}{arXiv:1301.2235}.

\bibitem{Ser96}
Serganova V., On generalizations of root systems, \href{https://doi.org/10.1080/00927879608825814}{\textit{Comm. Algebra}}
 \textbf{24} (1996), 4281--4299.

\bibitem{TV03}
Tarasov V., Varchenko A., Selberg-type integrals associated with
 {$\mathfrak{sl}_3$}, \href{https://doi.org/10.1023/B:MATH.0000010712.67685.9d}{\textit{Lett. Math. Phys.}} \textbf{65} (2003), 173--185,
 \href{https://arxiv.org/abs/math.QA/0302148}{arXiv:math.QA/0302148}.

\bibitem{TW13}
Tsuchiya A., Wood S., The tensor structure on the representation category of
 the {$\mathcal W_p$} triplet algebra, \href{https://doi.org/10.1088/1751-8113/46/44/445203}{\textit{J.~Phys.~A: Math. Theor.}}
 \textbf{46} (2013), 445203, 40~pages, \href{https://arxiv.org/abs/1201.0419}{arXiv:1201.0419}.

\bibitem{War09}
Warnaar S.O., A {S}elberg integral for the {L}ie algebra {$A_n$}, \href{https://doi.org/10.1007/s11511-009-0043-x}{\textit{Acta
 Math.}} \textbf{203} (2009), 269--304, \href{https://arxiv.org/abs/0708.1193}{arXiv:0708.1193}.

\bibitem{Wor89}
Woronowicz S.L., Differential calculus on compact matrix pseudogroups (quantum
 groups), \href{https://doi.org/10.1007/BF01221411}{\textit{Comm. Math. Phys.}} \textbf{122} (1989), 125--170.

\bibitem{Yam07}
Yamane H., Representations of a {$\mathbb Z/3\mathbb Z$}-quantum group,
 \href{https://doi.org/10.2977/PRIMS/1199403808}{\textit{Publ. Res. Inst. Math. Sci.}} \textbf{43} (2007), 75--93.

\end{thebibliography}
